*   Professional Mathematics teacher for Γ Λυκείου (12th grade).
    *   Task: Solve a specific LaTeX exercise about moving points $M$ and $N$.
    *   Constraint: Strict formatting guidelines (Thema environment, Erwthma environment, itemization, alignment, specific LaTeX commands).
    *   Specific exercise:
        *   $M$ starts at $O(0,0)$, moves on $Ox$ with $v_M = 3 \text{ m/s}$.
        *   $N$ starts at $y_0 = 10$ on $Oy$, moves away from origin with $v_N = 4 \text{ m/s}$.
        *   Q1: $M(t)$ and $N(t)$ coordinates.
        *   Q2: Distance $d(t) = MN$.
        *   Q3: Rate of change of $d(t)$ at $t=2$.
        *   Q4: Existence of minimum distance.

    *   $M$ starts at $(0,0)$, speed $3$ on $Ox$. So $x_M(t) = 3t$, $y_M(t) = 0$. $M(t) = (3t, 0)$.
    *   $N$ starts at $(0, 10)$, speed $4$ on $Oy$ (away from origin). So $x_N(t) = 0$, $y_N(t) = 10 + 4t$. $N(t) = (0, 10 + 4t)$.

    *   $d(t) = \sqrt{(x_N - x_M)^2 + (y_N - y_M)^2}$
    *   $d(t) = \sqrt{(0 - 3t)^2 + (10 + 4t - 0)^2}$
    *   $d(t) = \sqrt{9t^2 + (10 + 4t)^2}$
    *   $d(t) = \sqrt{9t^2 + 100 + 80t + 16t^2}$
    *   $d(t) = \sqrt{25t^2 + 80t + 100} = 5\sqrt{t^2 + 3.2t + 4}$.

    *   Rate of change is $d'(t)$.
    *   $d(t) = (25t^2 + 80t + 100)^{1/2}$.
    *   $d'(t) = \frac{1}{2}(25t^2 + 80t + 100)^{-1/2} \cdot (50t + 80)$.
    *   $d'(t) = \frac{25t + 40}{\sqrt{25t^2 + 80t + 100}}$.
    *   At $t=2$:
        *   $d(2) = \sqrt{25(4) + 80(2) + 100} = \sqrt{100 + 160 + 100} = \sqrt{360} = 6\sqrt{10}$.
        *   $d'(2) = \frac{25(2) + 40}{\sqrt{360}} = \frac{90}{6\sqrt{10}} = \frac{15}{\sqrt{10}} = \frac{15\sqrt{10}}{10} = 1.5\sqrt{10} \text{ m/s}$.

    *   To find the minimum, look for $d'(t) = 0$.
    *   $d'(t) = \frac{25t + 40}{\sqrt{25t^2 + 80t + 100}}$.
    *   $d'(t) = 0 \Rightarrow 25t + 40 = 0 \Rightarrow t = -40/25 = -1.6$.
    *   Since $t \ge 0$ (the movement starts at $t=0$), $d'(t)$ is always positive for $t \ge 0$ ($25t + 40 > 0$).
    *   Therefore, $d(t)$ is \textbf{γνησίως αύξουσα} for $t \in [0, +\infty)$.
    *   The minimum distance occurs at the starting point $t=0$.
    *   Minimum distance $d(0) = \sqrt{100} = 10 \text{ m}$.

    *   `\begin{thema}{Γ}`
    *   `\begin{erwthma}`
    *   `\item`
    *   `\begin{align*}` for calculations.
    *   `\textbf{...}` for emphasis.
    *   `\eng{...}` for technical terms (though not many here, maybe "minimum" or just standard Greek).
    *   `\[ ... \]` for key results.

    *   Check: $x_M = 3t, y_M = 0$. $x_N = 0, y_N = 10+4t$. Correct.
    *   Check: $d(t) = \sqrt{9t^2 + (10+4t)^2} = \sqrt{25t^2 + 80t + 100}$. Correct.
    *   Check: $d'(t) = \frac{50t + 80}{2\sqrt{25t^2 + 80t + 100}} = \frac{25t + 40}{\sqrt{25t^2 + 80t + 100}}$. Correct.
    *   Check: $t=2 \Rightarrow d'(2) = \frac{50+40}{\sqrt{100+160+100}} = \frac{90}{\sqrt{360}} = \frac{90}{6\sqrt{10}} = \frac{15}{\sqrt{10}} = 1.5\sqrt{10}$. Correct.
    *   Check: Min distance. $d'(t) > 0$ for $t \ge 0$. Minimum at $t=0$. Correct.

    *   The prompt asks for a professional math teacher's response in LaTeX.
    *   Ensure use of `\begin{thema}{Γ}` and `\begin{erwthma}`.
    *   Use `\textbf{γνησίως αύξουσα}`.
    *   Use `\begin{align*}` for $d'(t)$ computation.
    *   Units: m/s.

    *Self-Correction during drafting:* The problem says "Υπάρχει χρονική στιγμή κατά την οποία η απόσταση των δύο σημείων να γίνεται ελάχιστη;". Since $d(t)$ is strictly increasing for $t \ge 0$, the "minimum" is at the boundary $t=0$. I should state that for $t \ge 0$, the function is strictly increasing, so the minimum occurs at $t=0$.

    *Wait, let's re-read the prompt for "forbidden" items:*
    - No indefinite integrals. (N/A here)
    - No 2nd derivative test for extrema. (I'll use the sign of the 1st derivative). (N/A here)
    - No differential equations. (N/A)
    - No MVT for integrals. (N/A)

    *Let's double check the distance calculation again.*
    $M(t) = (3t, 0)$
    $N(t) = (0, 10 + 4t)$
    $d^2(t) = (3t)^2 + (10+4t)^2 = 9t^2 + 100 + 80t + 16t^2 = 25t^2 + 80t + 100$.
    $d(t) = \sqrt{25t^2 + 80t + 100}$. Correct.

    *Rate of change at $t=2$:*
    $d'(t) = \frac{50t + 80}{2\sqrt{25t^2 + 80t + 100}} = \frac{25t + 40}{\sqrt{25t^2 + 80t + 100}}$.
    $d'(2) = \frac{50 + 40}{\sqrt{100 + 160 + 100}} = \frac{90}{\sqrt{360}} = \frac{90}{6\sqrt{10}} = \frac{15}{\sqrt{10}} = \frac{3\sqrt{10}}{2} = 1,5\sqrt{10} \approx 4,74 \text{ m/s}$.

    *Minimum distance:*
    Since $t \ge 0$, $25t + 40 > 0$ and $\sqrt{25t^2 + 80t + 100} > 0$.
    So $d'(t) > 0$ for all $t \in [0, +\infty)$.
    The function $d(t)$ is \textbf{γνησίως αύξουσα} on $[0, +\infty)$.
    The minimum occurs at the left endpoint $t=0$.

    *Formatting check:*
    - `\begin{thema}{Γ}`
    - `\begin{erwthma}`
    - `\item`
    - `\begin{align*}`
    - `\[ ... \]`

    Let's write.
%=========== ΛΥΣΗ - {ΘΕΜΑ Γ} ===========

\begin{thema}{Γ}
\begin{erwthma}
\item Το σημείο $M$ κινείται στον ημιάξονα $Ox$ με σταθερή ταχύτητα $v_M = 3\si{m/s}$ ξεκινώντας από το $O(0,0)$, άρα η θέση του είναι:
\[ M(t) = (3t, 0) \]
Το σημείο $N$ κινείται στον ημιάξονα $Oy$ με σταθερή ταχύτητα $v_N = 4\si{m/s}$ ξεκινώντας από το $y_0 = 10$ και απομακρυνόμενο από την αρχή των αξόνων, άρα η θέση του είναι:
\[ N(t) = (0, 10 + 4t) \]
για $t \geq 0$.

\item Η απόσταση $d(t)$ μεταξύ των σημείων $M$ και $N$ δίνεται από τον τύπο της απόστασης δύο σημείων στο επίπεδο:
\begin{align*}
d(t) &= \sqrt{(x_N - x_M)^2 + (y_N - y_M)^2} \\
&= \sqrt{(0 - 3t)^2 + (10 + 4t - 0)^2} \\
&= \sqrt{9t^2 + (100 + 80t + 16t^2)} \\
&= \sqrt{25t^2 + 80t + 100}
\end{align*}

\item Ο ρυθμός μεταβολής της απόστασης $d(t)$ είναι η παράγωγος της συνάρτησης $d(t)$ ως προς τον χρόνο $t$. 
Έχοντας $d(t) = (25t^2 + 80t + 100)^{1/2}$, η παράγωγος είναι:
\begin{align*}
d'(t) &= \frac{1}{2}(25t^2 + 80t + 100)^{-1/2} \cdot (50t + 80) \\
&= \frac{50t + 80}{2\sqrt{25t^2 + 80t + 100}} \\
&= \frac{25t + 40}{\sqrt{25t^2 + 80t + 100}}
\end{align*}
Για $t = 2\si{s}$, ο ρυθμός μεταβολής είναι:
\[ d'(2) = \frac{25(2) + 40}{\sqrt{25(2)^2 + 80(2) + 100}} = \frac{50 + 40}{\sqrt{100 + 160 + 100}} = \frac{90}{\sqrt{360}} \]
\[ d'(2) = \frac{90}{6\sqrt{10}} = \frac{15}{\sqrt{10}} = \frac{15\sqrt{10}}{10} = 1,5\sqrt{10}\si{m/s} \]

\item Για να εξετάσουμε αν υπάρχει χρονική στιγμή ελάχιστης απόστασης, μελετάμε το πρόσημο της $d'(t)$ για $t \geq 0$:
\begin{itemize}
    \item Ο αριθμητής $25t + 40$ είναι πάντα θετικός για κάθε $t \geq 0$.
    \item Ο παρονομαστής $\sqrt{25t^2 + 80t + 100}$ είναι πάντα θετικός για κάθε $t \geq 0$.
\end{itemize}
Συνεπώς, $d'(t) > 0$ για κάθε $t \in [0, +\infty)$, γεγονός που σημαίνει ότι η συνάρτηση $d(t)$ είναι \textbf{γνησίως αύξουσα} στο διάστημα ορισμού της. 
Άρα, η ελάχιστη απόσταση επιτυγχάνεται στη χρονική στιγμή $t = 0$, όπου:
\[ d(0) = \sqrt{100} = 10\si{m} \]
\end{erwthma}
\end{thema}

%=========== ΤΕΛΟΣ ΛΥΣΗΣ - {ΘΕΜΑ Γ} ===========